\section{Related Work}
\label{sec:related-work}

\subsection{Metadata Extraction from Scientific Publications}

Metadata extraction from scientific publications is closely related to
document-level scientific information extraction. SciREX introduced
document-level extraction of salient entities and n-ary relations from
scientific articles \citep{jain2020scirex}, while SciER provides fine-grained
entity and relation annotations over full-text scientific publications
\citep{zhang2024scier}.

Related work has also considered the extraction of richer structured records
from scientific publications. Ghosh et al. study schema-based reconstruction
of curated materials-science datasets from full manuscripts
\citep{ghosh2024reliable}, while MatViX addresses multimodal structured
extraction across text, tables, and figures
\citep{khalighinejad2025matvix}. Closest to our setting, MOLE formulates
schema-driven metadata extraction from complete scientific documents
\citep{alyafeai2025mole}.

\subsection{Evaluation of Structured Extraction}

Structured extraction often requires different notions of correctness for
different fields. MOLE evaluates scientific metadata according to schema
answer type \citep{alyafeai2025mole}, while ExtractBench assigns field-level
metrics according to schema semantics \citep{ferguson2026extractbench}.

Repeated structured objects introduce an additional correspondence problem:
predicted and gold objects must be aligned before their contents can be
compared. Ghosh et al. use domain-specific criteria followed by greedy
matching to align extracted and gold rows in materials-science datasets
\citep{ghosh2024reliable}. MatViX instead aligns predicted and gold samples
through maximum bipartite matching before evaluating their properties
\citep{khalighinejad2025matvix}. ExtractBench uses LLM-based semantic alignment
for unordered object arrays and recursively evaluates the contents of matched
objects \citep{ferguson2026extractbench}.

Like prior structured-extraction evaluations, our evaluation must first align
predicted entities with gold entities. This is complicated by
publication-specific naming and entity boundaries, particularly for samples 
and datasets. We therefore use benchmark-specific alignment criteria
supported by evidence in the source PDFs.

\subsection{PaN Metadata and Research Data Management}

FAIR data management is an established priority in the photon and neutron
community, with metadata playing a central role in making experimental data
findable and reusable \citep{murphy2025fair}. NeXus provides
a widely used structured data format for neutron, X-ray, and muon experiments
\citep{konnecke2015nexus}, while the ExPaNDS project developed recommendations
for FAIR PaN data management and associated metadata practices
\citep{expands2022fair}.

PaNET provides a taxonomy of experimental techniques for photon and neutron
science \citep{collins2021panontologies}. Recent work proposes extending PaNET
with logical definitions to support automated reasoning and improve data
discovery \citep{tan2025panet}. DAPHNE4NFDI further provides recommendations
for metadata capture and specification spanning bibliographic/administrative and
experimental metadata \citep{lohstroh2025daphne}.

These efforts primarily address metadata representation, capture,
interoperability, and management across the experimental lifecycle. Our
setting considers the experimental context that can be recovered from scientific 
publications.
